\documentclass[12pt]{article}
\usepackage{geometry}
\geometry{a4paper, margin=1in}
\usepackage{hyperref}
\usepackage{booktabs}
\usepackage{enumitem}
\usepackage{xcolor}

\usepackage{xepersian}
\settextfont{Vazir}

\hypersetup{
    colorlinks=true,
    linkcolor=blue,
    filecolor=magenta,      
    urlcolor=cyan,
}

\begin{document}

\begin{center}
    \large\textbf{دانشکده علوم ریاضی - دانشگاه فردوسی مشهد} \\
    \large\textbf{سوالات و راهنمای ارزیابی جلسه دفاع پروژه پایگاه داده} \\
    \textbf{طراحی پایگاه داده سامانه مسیریابی شهری} \\
    \textbf{استاد درس: دکتر جلال‌الدین نصیری} --- \textbf{بهار ۱۴۰۵}
\end{center}

\hrule
\vspace{0.4cm}

این سند به عنوان راهنمای ارزیابی در جلسه ارائه و دفاع پروژه طراحی پایگاه داده سامانه مسیریابی شهری تنظیم شده است. سوالات به تفکیک فازهای پروژه و اهداف ارزیابی دسته‌بندی شده‌اند تا بتوانید میزان تسلط هر یک از اعضای گروه‌های دو نفره و اصالت کار آن‌ها را بسنجید.

\section{ارزیابی عمومی و تفکیک کار تیمی (Individual Contribution)}
با توجه به دو نفره بودن گروه‌ها، ابتدا باید مطمئن شوید هر دو دانشجو در پروژه نقش فعال داشته‌اند.
\begin{itemize}[label=$\bullet$, leftmargin=*]
    \item \textbf{سوال ۱:} «هر کدام از شما دقیقاً چه بخشی از مدل‌سازی ER و پیاده‌سازی کوئری‌ها را انجام داده‌اید؟ لطفاً یکی از جدول‌ها/کوئری‌هایی که خودتان طراحی کرده‌اید را نام برده و منطق آن را توضیح دهید.»
    \item \textbf{سوال ۲:} «اگر بخواهید ساختار فعلی پایگاه داده را برای پشتیبانی از ویژگی جدیدی مثل "سیستم امتیازدهی کاربران به وضعیت ترافیک مسیرها" توسعه دهید، چه جدول(ها) یا رابطه‌هایی را و به چه شکل به سیستم اضافه می‌کنید؟» (ارزیابی تفکر لحظه‌ای و طراحی پویا)
\end{itemize}

\section{سوالات تخصصی فاز اول: طراحی مدل ER (مفهومی)}
در این فاز تمرکز بر نحوه مدل‌سازی پدیده‌های دنیای واقعی (محدودیت‌ها، رویدادها، ترافیک و ایستگاه‌ها) در قالب موجودیت‌ها و روابط است.
\begin{itemize}[label=$\bullet$, leftmargin=*]
    \item \textbf{سوال ۳ (مدل‌سازی محدودیت‌های گردش - Turn Restrictions):} \\
    «چگونه محدودیت‌های گردش (مانند ممنوعیت گردش به چپ در یک تقاطع از مسیر A به مسیر B) را در مدل ER نشان داده‌اید؟ آیا تقاطع (نقطه) را به عنوان یک رابطه مدل کرده‌اید یا موجودیت؟ چرا؟»
    \item \textbf{سوال ۴ (پشتیبانی از حمل‌ونقل ترکیبی):} \\
    «مدل‌سازی ایستگاه‌های حمل‌ونقل عمومی و ترتیب آن‌ها در یک خط (مثلاً خط مترو یا اتوبوس) چگونه انجام شده است؟ چگونه تضمین کرده‌اید که یک ایستگاه می‌تواند در چند خط مختلف با ترتیب‌های متفاوت قرار بگیرد؟»
    \item \textbf{سوال ۵ (موجودیت‌های ضعیف - Weak Entities):} \\
    «آیا در طراحی خود از موجودیت ضعیف استفاده کرده‌اید؟ (مثلاً برای رویدادهای غیرعادی مسیرها یا بازه‌های زمانی خاص). علت نیاز به موجودیت ضعیف در آن بخش چه بود و کلید شناسگر (Partial Key) آن چیست؟»
    \item \textbf{سوال ۶ (روابط چندبه‌چند):} \\
    «رابطه بین وسایل نقلیه و مسیرها را چگونه مدل کرده‌اید؟ با توجه به اینکه هر مسیر قوانین خاصی برای هر نوع وسیله نقلیه دارد، ویژگی‌های این قوانین (مانند حداکثر سرعت مجاز یا ساعت ممنوعیت عبور) را در کدام موجودیت یا رابطه قرار داده‌اید؟»
\end{itemize}

\newpage

\section{سوالات تخصصی فاز دوم: تبدیل به مدل رابطه‌ای و نرمال‌سازی}
هدف این بخش سنجش درک دانشجو از مفاهیم وابستگی‌های تابعی (FD)، نرمال‌سازی و سناریوهای غیرنرمال‌سازی (Denormalization) است.
\begin{itemize}[label=$\bullet$, leftmargin=*]
    \item \textbf{سوال ۷ (شناسایی وابستگی‌های تابعی):} \\
    «یکی از کلیدی‌ترین وابستگی‌های تابعی (FD) غیربدیهی که در جدول مسیرها (Routes) یا قوانین زمانی شناسایی کردید را بنویسید. این وابستگی چگونه روی کلید کاندید جدول تاثیر گذاشت؟»
    \item \textbf{سوال ۸ (تحلیل سطح نرمال‌سازی):} \\
    «جداول نهایی شما در چه سطح نرمال‌سازی (مثلاً 3NF یا BCNF) قرار دارند؟ روند نرمال‌سازی یکی از جدول‌های پیچیده را گام‌به‌گام توضیح دهید.»
    \item \textbf{سوال ۹ (توجیه غیرنرمال‌سازی - Denormalization):} \\
    «در صورت استفاده از Denormalization در داده‌های ترافیکی یا زمانی، چه توجیهی برای آن دارید؟ آیا بهبود سرعت جستجو ارزش افزایش احتمال بروز ناهنجاری (Anomaly) را در این بخش داشت؟ چگونه ناهنجاری‌ها را کنترل می‌کنید؟»
    \item \textbf{سوال ۱۰ (جامعیت ارجاعی - Referential Integrity):} \\
    «در صورت حذف یک مسیر (Route)، وضعیت رویدادهای ثبت‌شده روی آن مسیر یا محدودیت‌های گردش متصل به آن چه می‌شود؟ قواعد \lr{ON DELETE} را برای کلیدهای خارجی مربوطه چگونه تنظیم کرده‌اید؟ (\lr{CASCADE}، \lr{SET NULL} یا \lr{RESTRICT})»
\end{itemize}

\section{سوالات تخصصی فاز سوم: جبر رابطه‌ای و کوئری‌های SQL}
این سوالات میزان تسلط بر ترجمه منطق جبر رابطه‌ای به SQL و استفاده از مفاهیم پیشرفته SQL را ارزیابی می‌کنند.
\begin{itemize}[label=$\bullet$, leftmargin=*]
    \item \textbf{سوال ۱۱ (جبر رابطه‌ای در مقابل SQL):} \\
    «در سوال ۴ جبر رابطه‌ای (پیدا کردن مسیرهای متصل در تقاطع مشترک که گردش بین آن‌ها ممنوع است)، نحوه اعمال عملگرهای پیوند (Join) و تفاضل (Difference) را روی کاغذ بنویسید. پیاده‌سازی این منطق در SQL به چه شکل انجام شد؟»
    \item \textbf{سوال ۱۲ (توانمندی توابع پنجره‌ای - Window Functions):} \\
    «در سوال ۳ بخش SQL، برای پیدا کردن نقاط با مسیرهای متصل بیشتر از میانگین، چطور از Window Function استفاده کردید؟ تفاوت استفاده از \lr{OVER()} با یک \lr{GROUP BY} معمولی در این سناریو چیست؟»
    \item \textbf{سوال ۱۳ (کوئری‌های بازگشتی - Recursive CTE):} \\
    «در سوال ۵ بخش SQL (یافتن نقاط قابل دسترس با حداکثر ۳ واسطه)، بخش پایه (Anchor Member) و بخش بازگشتی (Recursive Member) در کوئری شما چیست؟ شرط توقف این بازگشت را چگونه تعریف کردید تا در حلقه‌های بی‌نهایت (گراف‌های حلقوی شهر) گرفتار نشود؟»
    \item \textbf{سوال ۱۴ (کارایی و ایندکس‌گذاری):} \\
    «با توجه به کوئری مقایسه ترافیک در ساعات اوج (سوال ۴ SQL)، اجرای این کوئری روی داده‌های بزرگ بسیار سنگین خواهد بود. برای بهینه‌سازی آن چه ایندکس‌هایی (Indexes) روی جدول داده‌های زمانی یا ترافیکی پیشنهاد می‌کنید؟»
\end{itemize}

\section{سوالات تخصصی فاز چهارم (امتیازی): تحلیل زمانی و رتبه‌بندی}
بخش‌های امتیازی معمولاً خلاقیت و درک عمیق‌تر سیستم‌های واقعی را نشان می‌دهند.
\begin{itemize}[label=$\bullet$, leftmargin=*]
    \item \textbf{سوال ۱۵ (طراحی View پویا و مدیریت زمان):} \\
    «در سوال اول فاز امتیازی، View طراحی شده چگونه فیلتر کردن مسیرهای مجاز در "زمان جاری سیستم" را انجام می‌دهد؟ اگر فیلدهای \lr{valid\_from} و \lr{valid\_to} شامل مقادیر \lr{NULL} باشند (به معنی نامحدود بودن زمان اعتبار)، شرط‌های بررسی زمان در View شما چطور این موضوع را هندل می‌کنند؟»
    \item \textbf{سوال ۱۶ (تفاوت RANK و ROW\_NUMBER):} \\
    «در سوال دوم امتیازی، برای رتبه‌بندی سریع‌ترین مسیرها بر حسب نوع وسیله نقلیه، تفاوت خروجی کوئری شما در صورت استفاده از \lr{RANK()} به جای \lr{ROW\_NUMBER()} چیست؟ در چه حالتی رفتار این دو تابع متفاوت خواهد بود؟»
\end{itemize}

\newpage

\section{راهنمای بارم‌بندی پیشنهادی در جلسه ارائه}
\begin{table}[h!]
\centering
\caption{بارم‌بندی پیشنهادی برای بخش‌های مختلف ارائه}
\vspace{0.3cm}
\begin{tabular}{lcl}
\toprule
\textbf{بخش ارزیابی} & \textbf{درصد نمره} & \textbf{معیار ارزیابی} \\
\midrule
صحت و اصالت ERD & ۲۵٪ & توانایی توضیح روابط چندبه‌چند، موجودیت‌های ضعیف و محدودیت‌ها \\
تحلیل نرمال‌سازی & ۲۰٪ & ارائه مستندات شفاف از FDs و توجیه مناسب برای ساختار جداول \\
تسلط بر کوئری‌نویسی & ۳۵٪ & نوشتن یا اصلاح زنده یک کوئری (مثلاً تغییر فیلتر یک Recursive CTE) \\
پاسخ به چالش‌ها و امتیازها & ۲۰٪ & نحوه مدیریت بازه‌های زمانی در Viewها و تسلط بر توابع رتبه‌بندی \\
\bottomrule
\end{tabular}
\end{table}

\end{document}
